% 自编教学补充；阅读来源见本章 README。
\begin{frame}[fragile]{5.2A 赋值是起别名，复制才得到新列表}
\small 变量是指向对象的名字。赋值不会自动复制列表；copy() 创建一个新的外层列表。先预测三个列表的内容，再执行。
\medskip
\begin{lstlisting}
original = [80, 90]
alias = original
copied = original.copy()
alias.append(100)
print(original)          # [80, 90, 100]
print(copied)            # [80, 90]
print(alias is original) # True
\end{lstlisting}
\end{frame}
\begin{frame}[fragile]{课堂练习：5.2A 赋值是起别名，复制才得到新列表}
把 alias[0] 改成 60，哪些列表会改变？如果列表元素本身还是列表，copy() 会复制内层吗？
\medskip
\begin{block}{检查思路}
先写出预期结果，再运行。参考答案见配套 Notebook 的折叠单元格。
\end{block}
\end{frame}

\begin{frame}[fragile]{5.2B sort 与 sorted：先看返回值}
\small sort() 修改原列表并返回 None；sorted() 返回新列表。不能把 names.sort() 的返回值当作排好序的名单。
\medskip
\begin{lstlisting}
names = ["Wang", "Li", "Zhang"]
ordered = sorted(names, key=len)
print(names, ordered)
result = names.sort()
print(names)
print(result)  # None
\end{lstlisting}
\end{frame}
